\chapter{系统测试与分析}

本章详细记录了系统在开发完成后的功能测试过程及结果分析。测试采用黑盒测试方法，通过模拟科研人员与管理员的实际操作，验证各功能模块是否达到设计要求与业务规范。

\section{测试环境与测试策略}
本系统的测试环境采用容器化集群部署，后端基于 Spring Boot 框架，数据库选用 MySQL 8.0。测试工作重点围绕报销审计流程的准确性、预警算法的有效性以及系统权限的安全性展开，旨在确保科研经费管理系统的稳定运行与数据严谨性。

\section{核心功能测试用例}
本轮测试共设计了 6 组具有代表性的测试用例，主要涵盖了系统认证、预算校验逻辑以及智能预警模块等关键功能区域。

\subsection{系统认证与业务逻辑测试}
系统对基础认证与业务逻辑异常操作的测试反馈如表~\ref{tab:test_basic}~所示：

\begin{table}[htbp]
    \centering
    \caption{系统基础认证与财务逻辑测试表}
    \label{tab:test_basic}
    \zihao{-4}\songti
    \begin{tabularx}{\textwidth}{ccYYYc}
    \toprule
    \textbf{序号} & \textbf{测试项名称} & \textbf{输入数据} & \textbf{预期结果} & \textbf{实际结果} & \textbf{结论} \\ \midrule
    01 & 登录安全性验证 & 用户输入错误密码 & 系统拒绝访问并提示错误 & 成功拦截非法登录请求 & 通过 \\
    02 & 预算总额校验 & 预算分项合计超过总额 & 系统提示分配金额超出限制 & 成功阻止保存并提示错误 & 通过 \\
    03 & 超额报销拦截 & 科目余额不足时提交报销 & 系统提示科目余额不足 & 报销流程中止并提示原因 & 通过 \\ \bottomrule
    \end{tabularx}
\end{table}

\subsection{智能预警效果与安全性测试}
该环节重点验证系统在多维度数据分析及权限访问控制方面的可靠性，测试结果如表~\ref{tab:test_smart}~所示：

\begin{table}[htbp]
    \centering
    \caption{智能决策预警与安全性测试表}
    \label{tab:test_smart}
    \zihao{-4}\songti
    \begin{tabularx}{\textwidth}{ccYYYc}
    \toprule
    \textbf{序号} & \textbf{测试项名称} & \textbf{输入数据} & \textbf{预期结果} & \textbf{实际结果} & \textbf{结论} \\ \midrule
    01 & 偏差识别预警测试 & 执行进度明显滞后于时间进度 & 自动生成“执行滞后”标签 & 预警状态触发正常 & 通过 \\
    02 & 垂直权限控制测试 & 普通用户尝试请求管理接口 & 返回 403 拒绝访问错误 & 后端成功拦截越权请求 & 通过 \\
    03 & 数据一致性验证 & 核准一笔报销申请 & 相关统计看板即时更新 & 各维度数据同步更新正常 & 通过 \\ \bottomrule
    \end{tabularx}
\end{table}

\section{测试分析}
根据上述测试用例的执行结果，系统各项功能均达到了预期设计目标。系统的逻辑严密性得到了有效验证，特别是内置的预算校验与财务预警机制，能够有效防止违规报销，确保了经费管理的规范性。此外，系统通过定时扫描任务，能够准确识别项目执行进度的滞后情况，为科研管理部门提供了有效的决策支持。

\section{本章小结}
本章对科研经费管理系统的核心功能、权限安全及智能预警模块进行了全面的测试。测试结果显示，系统运行稳定，各项拦截逻辑和预警算法均发挥了预期作用。通过测试验证，本系统不仅提升了经费管理的自动化水平，还增强了财务监管的精准度，具备良好的实际应用价值。
